\section{Acceptance Probability Algorithm}
\label{sec:acceptance_probability_algorithm}
In this section,
$\mathbb{F}_2$ denotes the finite field with $2$ elements.
For $\v v \in \mathbb{F}_2^n$,
write $P_{\v v} :=\bigotimes_{q =1}^n P^{v_q}$ for $P \in\{\X, \Z\}$.

\begin{prop}
\label{prop:probability_runtime}
\Cref{alg:probability_trivial_syndrome} computes
\cref{eq:probability_trivial_syndrome}
in $\mathcal O(n^3+N_\alpha k^2)$ time,
where $N_\alpha \le 4^k$ is the number of nonzero coefficients
$\alpha_L$ defining $\rho$ by \cref{eq:pauli_decomposition}.
\end{prop}

\begin{algorithm}[H]
\caption{Probability of a trivial syndrome after suffering a Clifford error.}
\label{alg:probability_trivial_syndrome}
\begin{algorithmic}[1]
\Require Encoder $C$, Clifford error $E$,
and logical coefficients $\{\alpha_L:L\in\mathcal P_k\}$.
\Ensure $\pr(\v s=\v 0)$.
\State Set $F\gets C^\dag EC$ and build $R$ as in \cref{eq:forget_check_Z_matrix}.
\State Row-reduce $R$ to obtain $\rk R$, a basis $\mathcal{B}$ for $\ker R$,
an image test, and a solver for $R\v x=\v y$.
\State Precompute the data to evaluate $\chi$ using
\cref{eq:chi_phase_formula}.
\ForAll{$\v e \in \mathcal{B}$}
	\If{$\chi(\v e)=1$}
		\State \Return $0$
	\EndIf
\EndFor
\State $\eta\gets0$
\ForAll{$L\in\mathcal P_k$ with $\alpha_L\ne0$}
	\If{$\v \rho(\I^{\otimes r}\otimes L) \in\im R$}
	\label{line:image_membership_test}
		\State $\eta\gets\eta+(-1)^{\chi(\v u_L)}\alpha_L$,
		where $\v u_L$ is any solution to
		$R\v u_L =\v \rho(\I^{\otimes r} \otimes L)$
	\EndIf
\EndFor
\State \Return $2^{-\rk R}\eta$
\end{algorithmic}
\end{algorithm}

\begin{proof}
Let $r :=n-k$.
It helps to work in unencoded space:
choose an encoder $C$ such that the encoded state is
$\overline{\rho} =C \big(\proj{0}^{\otimes r} \otimes \rho\big) C^\dag$.
Then the unencoded stabiliser generators are simply
$\mathcal Z:=\{\Z_1,\ldots,\Z_r\}$
and the unencoded Clifford error is $F:=C^\dag EC$.
Write the transformed stabiliser element corresponding to
$\v u\in\mathbb{F}_2^r$ as $\Z'_{\v u}:=F^\dag\Z_{\v u}F$.
In unencoded space, \cref{eq:trace_triple} becomes
\begin{equation}
\pr(\v s=\v 0)
=
2^{-2r-k}
\sum_{L\in\mathcal P_k}\alpha_L
\sum_{\v u,\v v\in\mathbb{F}_2^r}
\tr\left[\Z_{\v u}'\Z_{\v v}(\I^{\otimes r}\otimes L)\right].
\label{eq:single_conjugated_trace_expansion}
\end{equation}
As before,
we need only consider products
$\Z_{\v u}'\Z_{\v v}(\I^{\otimes r}\otimes L)$
that are proportional to $\I$.
For a given $L$,
this occurs if there exists a $\Z_{\v u}'$ whose...
\begin{enumerate}
	\item first $r$ tensor factors are in $\{\I, \Z\}$,
	\item last $k$ tensor factors match $L$ up to a sign.
\end{enumerate}
These two conditions can be restated as
the single equation
$R\v u=\v \rho(\I^{\otimes r}\otimes L)$,
where we have defined
$\v \rho(P)
:=(\v x^\T\mid z_{r+1},\ldots,z_n)^\T
\in\mathbb{F}_2^{n+k}$
for a Hermitian $n$-qubit Pauli $P$
whose binary symplectic vector is
$(\v x^\T \mid\v z^\T)^\T$,
and the matrix
\begin{equation}
R:=\big[\v \rho(F^\dag\Z_1F),\ldots,\v \rho(F^\dag\Z_rF)\big].
\label{eq:forget_check_Z_matrix}
\end{equation}
There can be multiple solutions to
$R\v u=\v \rho(\I^{\otimes r}\otimes L)$;
we define the \emph{fibre} over $L$ as the solution set:
\begin{equation}
\Phi_L:=
\{\v u\in\mathbb{F}_2^r:R\v u=\v \rho(\I^{\otimes r}\otimes L)\}.
\label{eq:fibre_defn}
\end{equation}
For $\v u\in\Phi_L$,
there is a unique $\v v(\v u)\in\mathbb{F}_2^r$ such that
$\Z_{\v v(\v u)}$ matches the first $r$ tensor factors of $\Z_{\v u}'$
i.e.
\begin{equation}
\Z'_{\v u}\Z_{\v v(\v u)}(\I^{\otimes r}\otimes L)
=(-1)^{\chi(\v u)}\I.
\label{eq:chi_defn}
\end{equation}
\Cref{lem:trivial_syndrome_phase_evaluation}
gives an explicit formula for $\chi: \mathbb{F}_2^r \to \mathbb{F}_2$
and shows that it depends only on $\v u$, not on $L$.

The zero case of \cref{theorem:probability_trivial_syndrome}
is detected on $\ker R$:
by \cref{eq:chi_defn} with $L=\I$,
$\chi(\v k)=1$ iff
$\Z'_{\v k}\in-\langle\mathcal Z\rangle$.
By \cref{lem:kernel_chi_linear},
checking $\chi$ on a basis of $\ker R$
detects precisely whether this happens.

Assume the algorithm does not return zero.
\Cref{line:image_membership_test} selects
exactly the $L$'s satisfying \cref{eq:logical_membership}.
For any solution $\v u_L\in\Phi_L$,
\cref{eq:chi_defn} identifies the sign in
\cref{eq:probability_trivial_syndrome} as
$\omega_L=(-1)^{\chi(\v u_L)}$.
This sign is solution-independent by
\cref{theorem:probability_trivial_syndrome}.
Finally,
$\rk R=s$
because $R$ records the transformed stabilisers modulo
the original unencoded stabiliser group.
Thus the accumulator $\eta$ is the sum in
\cref{eq:probability_trivial_syndrome},
and the algorithm returns the prefactor $2^{-s}=2^{-\rk R}$.

It remains to prove the stated time complexity:
\begin{enumerate}
	\item Tableau composition computes
	$F=C^\dag EC$ in $\mathcal O(n^3)$ time
	\cite[\S III]{Aaronson2004}
	\cite[\S 2.3.2]{Gidney2021a},
	after which the columns $F^\dag\Z_iF$ are read from the tableau.
	The same preprocessing row-reduces $R$, builds the image-membership test
	and precomputes the phase data in
	\cref{eq:chi_phase_formula}; checking $\chi$ on a basis of $\ker R$
	is also $\mathcal O(n^3)$ time.
	\item For each nonzero coefficient $\alpha_L$:
	one need not explicitly form $\v u_L\in\mathbb{F}_2^r$.
	After row-reduction, the membership test and
	$\chi(\v u_L)$ can be precomposed with the binary symplectic
	vector of $L$. This gives linear and quadratic forms in $2k$ variables,
	respectively, which can be evaluated in $\mathcal O(k^2)$ time.
	\qedhere
\end{enumerate}
\end{proof}

\noindent
The following kernel fact is used in the proof.

\begin{lem}
\label{lem:kernel_chi_linear}
The restriction $\chi|_{\ker R}$ is linear
i.e.\ $\chi(\v k+\v l)
=\chi(\v k)+\chi(\v l)\pmod 2$ for $\v k, \v l \in \ker R$.
\end{lem}
\begin{proof}
$
(-1)^{\chi(\v k+\v l)}\I
\overset{\eqref{eq:chi_defn}}=
\Z_{\v k+\v l}'\Z_{\v v(\v k+\v l)}
=
\Z_{\v k}'\Z_{\v v(\v k)}
\Z_{\v l}'\Z_{\v v(\v l)}
\overset{\eqref{eq:chi_defn}}=
(-1)^{\chi(\v k)+\chi(\v l)}
\I$.
\end{proof}

\begin{lem}
\label{lem:trivial_syndrome_phase_evaluation}
Let
\[
F^\dag\Z_iF=\i^{\phi_i}\X_{\v x_i}\Z_{\v z_i},
\qquad
\v \phi=(\phi_1,\ldots,\phi_r)^\T\in\mathbb Z_4^r,
\qquad
X=[\v x_1,\ldots,\v x_r],
\qquad
Z=[\v z_1,\ldots,\v z_r].
\]
Then the function $\chi$ in \cref{eq:chi_defn} is
\begin{equation}
\chi(\v u)
:=
\frac{\v \phi^\T\v u-(X\v u)^\T(Z\v u)}2
+\sum_{i<j} u_i u_j\,\v z_i^\T\v x_j
\pmod 2,
\label{eq:chi_phase_formula}
\end{equation}
where the matrix--vector products $X\v u$ and $Z\v u$ are over $\mathbb{F}_2$,
then interpreted as $0/1$ integer vectors in the numerator.
The half is well-defined on this domain because
the numerator is always even.
\end{lem}
\begin{proof}
Consider first
\begin{align}
\Z_{\v u}'
&=\prod_{j \in \supp \v u} F^\dag \Z_j F \notag \\
&=\i^{\v \phi^\T \v u} \prod_{j \in \supp \v u}
\X_{\v x_j} \Z_{\v z_j}
\end{align}
We want to express this product
as one pure-X Pauli followed by one pure-Z Pauli.
In moving
all X components to the left and
all Z components to the right,
we accumulate factors of $-1$
because Paulis either commute or anticommute.
Specifically,
each time we move an $\X_{\v x_j}$
past a $\Z_{\v z_i}$ for $i <j$,
we pick up a $(-1)^{\v z_i^\T \v x_j}$ factor.
Thus,
the total number of $-1$ factors accumulated is
$\sum_{i <j} u_i u_j \v z_i^\T \v x_j$.
The X-component vector is $\sum_{j \in \supp \v u} \v x_j =X\v u$;
similarly for Z.
So the result is
\begin{equation}
\label{eq:raw_transformed_unencoded_product}
\Z'_{\v u}
=
\i^{\v \phi^\T\v u}
(-1)^{\sum_{i<j} u_i u_j\,\v z_i^\T\v x_j}
\X_{X\v u}\Z_{Z\v u}.
\end{equation}
Substituting this into \cref{eq:chi_defn},
we see that $\X_{X\v u}\Z_{Z\v u}$
equals $\Z_{\v v(\v u)}(\I^{\otimes r}\otimes L)$,
but without the powers of $\i$
from any $\Y$ tensor factors that appear in
$\Z_{\v v(\v u)}(\I^{\otimes r}\otimes L)$.
There are $(X\v u)^\T(Z\v u)$ such tensor factors
i.e.\ $\i^{(X\v u)^\T(Z\v u)} \X_{X\v u}\Z_{Z\v u}
=\Z_{\v v(\v u)}(\I^{\otimes r}\otimes L)$.
Substituting this into \cref{eq:raw_transformed_unencoded_product}
then into \cref{eq:chi_defn} implies the result.
\end{proof}